\section{Estado del Arte / Benchmarking}
Este capítulo revisa la literatura científica reciente centrada específicamente en la \textbf{aplicación} de técnicas de aprendizaje al pronóstico de generación fotovoltaica. A diferencia del marco teórico —que expuso la genealogía de los métodos—, aquí el interés recae en cómo esos métodos se han instrumentado sobre datos solares reales, qué desempeño empírico han reportado y qué limitaciones persisten. La revisión avanza de lo general a lo específico: parte de los enfoques de \textit{Machine Learning} y aprendizaje profundo aplicados al pronóstico solar, transita hacia las arquitecturas atencionales, y converge en la frontera que motiva esta investigación: la transferencia de conocimiento entre plantas bajo escasez de datos.

\subsection{Enfoques de Machine Learning y Aprendizaje Profundo en el Pronóstico Fotovoltaico}
La investigación aplicada inicial del pronóstico fotovoltaico se apoyó predominantemente en algoritmos ajustados con el historial climático de una sola ubicación. En una evaluación comparativa sobre un conjunto de datos solar real, Arslan et al. (2025) contrastaron algoritmos clásicos de \textit{Machine Learning} frente a modelos de aprendizaje profundo basados en Transformers, documentando la superioridad de estos últimos para capturar la dinámica no lineal de la irradiancia, aunque a costa de un mayor requerimiento computacional y de datos \cite{Arslan2025}. En la misma línea, los estudios centrados en algoritmos de potenciación del gradiente han consolidado a XGBoost como una línea base robusta y explicable para el pronóstico energético, destacando su eficiencia y su capacidad para jerarquizar la importancia de las variables meteorológicas, si bien evidencian su fragilidad ante regímenes de datos no observados en entrenamiento \cite{rodriguez2023energy}. Estos trabajos establecen el punto de partida empírico: los modelos locales son eficaces sobre bancos de datos extensos, pero su desempeño se degrada precisamente en las condiciones de escasez que caracterizan la puesta en marcha de una planta.

\subsection{Arquitecturas Atencionales Aplicadas a la Energía Solar}
La adopción de arquitecturas Transformer en el dominio solar ha producido resultados consistentemente favorables. López Santos et al. (2022) aplicaron el \textit{Temporal Fusion Transformer} al pronóstico fotovoltaico a un día vista (\textit{day-ahead}), aprovechando su capacidad para fusionar covariables estáticas y dinámicas y para emitir predicciones por cuantiles, y reportaron mejoras significativas frente a las líneas base recurrentes \cite{LopezSantos2022}. El-Saieed (2025) propuso un Transformer adaptativo para el pronóstico de largo plazo de temperatura y radiación en el contexto de la generación fotovoltaica, mostrando la idoneidad del mecanismo de atención para modelar las dependencias meteorológicas de largo alcance que gobiernan la producción \cite{ElSaieed2025}. Ampliando el espectro de fuentes de datos, Li et al. (2024) abordaron el \textit{nowcasting} solar mediante un Transformer que ingiere imágenes de satélite geoestacionario, evidenciando que la atención puede explotar información espacial exógena para anticipar la variabilidad inducida por la nubosidad \cite{Li2024}. En conjunto, esta línea confirma que las arquitecturas atencionales constituyen el estado del arte metodológico en el pronóstico fotovoltaico aplicado, pero la mayoría de estos trabajos opera todavía bajo el supuesto de disponer de historial local suficiente para el entrenamiento.

\subsection{Transferencia de Conocimiento y Escasez de Datos en Sistemas Fotovoltaicos}
La frontera de la disciplina, y el núcleo de esta investigación, es la aplicación de la transferencia de conocimiento para resolver la escasez de datos. Bak et al. (2025) testearon empíricamente esquemas de \textit{Transfer Learning} para el pronóstico de potencia fotovoltaica \textbf{entre regiones}, agregando datos de gran escala de múltiples plantas operativas para pre-entrenar un modelo base; sus resultados demuestran una atenuación de los márgenes de error durante las primeras semanas de operación en sitios geográficos nuevos, aunque el esquema aún requiere una fase de ajuste fino (\textit{fine-tuning}) una vez disponibles los primeros registros locales \cite{Bak2025}. Llevando esta idea a su expresión más extrema, Mishra et al. (2025) exploraron el pronóstico de irradiancia solar de corto plazo mediante modelos fundacionales en un régimen de \textbf{transferencia zero-shot}, es decir, sin ajuste alguno sobre la planta objetivo, anticipando el potencial de los modelos pre-entrenados a gran escala para operar desde el primer instante \cite{Mishra2025}. Estos trabajos validan la dirección del aprendizaje multi-sitio, pero dejan abierta la pregunta de cuánto de esa ganancia se sostiene bajo un protocolo de evaluación estrictamente ciego a la planta objetivo y contrastado de forma justa contra líneas base locales maduras.

\subsection{Síntesis Comparativa y Brecha Identificada}
La Tabla \ref{tab:benchmarking} sintetiza los estudios revisados según su enfoque, aporte principal y limitación metodológica.

\begin{table}[htbp]
  \centering
  \resizebox{\textwidth}{!}{
  \begin{tabular}{|m{3.2cm}|m{3.3cm}|m{4.5cm}|m{4.5cm}|}
    \hline
    \textbf{Estudio} & \textbf{Enfoque} & \textbf{Aporte principal} & \textbf{Limitación metodológica} \\
    \hline
    Arslan et al. (2025) & ML vs.\ Transformer (local) & Comparación empírica sobre datos solares reales; superioridad de la atención. & Entrenamiento local; no aborda el arranque en frío. \\
    \hline
    Zhu et al.\ (XGBoost) & Machine Learning local & Línea base robusta y explicable para pronóstico energético. & Degradación severa fuera del rango histórico observado. \\
    \hline
    López Santos et al. (2022) & TFT para FV (\textit{day-ahead}) & Fusión de covariables estáticas/dinámicas y pronóstico por cuantiles. & Requiere historial local del sitio evaluado. \\
    \hline
    El-Saieed (2025) & Transformer adaptativo FV & Modelado de dependencias meteorológicas de largo plazo. & Enfocado en variables meteorológicas, no en transferencia inter-sitio. \\
    \hline
    Li et al. (2024) & Transformer + satélite & Explotación de información espacial exógena (nubosidad). & Dependencia de infraestructura de datos satelitales. \\
    \hline
    Bak et al. (2025) & Transfer Learning inter-regional & Atenuación del error en las primeras semanas en sitios nuevos. & Exige \textit{fine-tuning} con datos locales iniciales. \\
    \hline
    Mishra et al. (2025) & Zero-shot / modelo fundacional & Pronóstico de irradiancia sin ajuste sobre el sitio objetivo. & Foco en irradiancia; sin \textit{benchmarking} contra locales maduros. \\
    \hline
  \end{tabular}
  }
  \caption{\label{tab:benchmarking} Síntesis comparativa de los estudios de pronóstico fotovoltaico aplicado revisados.}
\end{table}

De la revisión emerge una brecha clara. Si bien la literatura valida tanto la eficacia de las arquitecturas atencionales sobre datos solares como la promesa del aprendizaje por transferencia para mitigar la escasez de datos, los trabajos existentes rara vez combinan tres condiciones simultáneamente: \textbf{(i)} una evaluación estrictamente ciega a la planta objetivo (\textit{zero-shot} real, sin ajuste ni contexto local), \textbf{(ii)} un contraste directo y justo contra líneas base locales maduras bajo un protocolo idéntico, y \textbf{(iii)} una caracterización de la incertidumbre mediante intervalos probabilísticos y un análisis de sensibilidad estacional. La mayoría, además, opera sobre pocas plantas o requiere fases de ajuste fino que difuminan la condición de arranque en frío puro.

Para abordar esta brecha, el presente proyecto de título implementa y evalúa un marco predictivo \textit{Cross-Site} soportado sobre el \textit{Temporal Fusion Transformer}, el \textit{Informer} y N-HiTS, bajo un protocolo de retención de una planta (\textit{Leave-One-Plant-Out}) que garantiza la integridad \textit{zero-shot}. La evaluación se estructura como un \textit{benchmarking} cerrado que contrasta estas arquitecturas globales contra líneas base estrictamente locales (XGBoost y las mismas redes profundas entrenadas de forma aislada), reporta métricas puntuales y probabilísticas, y estratifica los resultados por estación del año sobre un conjunto de \textbf{94 plantas fotovoltaicas} del Sistema Eléctrico Nacional chileno. Esta configuración permite dictaminar cuantitativamente en qué medida el conocimiento climático generalizado de la flota operativa mitiga los sesgos predictivos durante el arranque en frío de nuevos parques solares.
